\section{Technical estimates}
\label{sec:analysis}
The analysis separates boundary and interior steps. Boundary steps yield
decrease-based gradient and curvature estimates, while interior steps
contract the gradient. A joint potential then controls the adaptive scale.

For the estimates in this section, fix a pair $(L_0,L_1)$ satisfying
Assumption~\ref{ass:setting} and assume $\Delta>0$. Define
\begin{equation}\label{eq:Q}
 \begin{aligned}
 Q&=L_0+L_1^{3/2}\Delta,\\
 Q_0&=\sigma_0+L_0+L_1\norm{\nabla f(x_0)}+L_1^{3/2}\Delta,
 \end{aligned}
\end{equation}
so $Q_0\ge Q>0$, and set
\begin{equation}\label{eq:logs}
 \mathcal L=1+\lp\frac{Q_0}{Q}
                +\lp\frac{Q_0}{\sigma_0}
                +\lp\frac{\norm{\nabla f(x_0)}}{(\Delta^2Q)^{1/3}}.
\end{equation}
We write $\widetilde O(1)$ for additive logarithmic terms in problem
parameters and initial data, independent of $\epsilon,\delta$;
the explicit term proved below is $O(\mathcal L)$.
All constants in $O(\cdot)$ and all occurrences of $C$ are universal;
$C$ may change between inequalities.

\subsection{Nonzero steps and finite backtracking}

Write $r_k=\norm{s_k}$.

\begin{lemma}[Nonzero trial steps]\label{lem:nonzero}
Every trial has $h_k,R_k,r_k>0$.
\end{lemma}
\begin{proof}
By \eqref{eq:model-scales}, $h_k=0$ would imply
$\nabla f(x_k)=0$ and $\nabla^2f(x_k)\succeq0$, so the stopping
test would have succeeded. Thus $h_k,R_k>0$.
If $\nabla f(x_k)\ne0$, \eqref{eq:kkt-stationarity}
rules out $s_k=0$. Otherwise failure of the stopping test gives
$\lambda_{\min}(\nabla^2f(x_k))<-\delta$, and
$2\sqrt{\sigma_kh_k}=-\tfrac14\lambda_{\min}(\nabla^2f(x_k))$.
The model Hessian therefore has a negative eigenvalue. A sufficiently
small nonzero step along a corresponding eigenvector has negative
model value, whereas the value at zero is zero. Hence $r_k>0$.
\end{proof}

\begin{lemma}[Sufficient condition for acceptance]\label{lem:safe}
A trial is accepted whenever
\begin{equation}\label{eq:safe}
 \sigma_k\ge L_0+L_1\norm{\nabla f(x_k)}
        +\sqrt{L_1}[-\lambda_{\min}(\nabla^2 f(x_k))]_+.
\end{equation}
Consequently, backtracking terminates after finitely many rejections.
\end{lemma}
\begin{proof}
Fix a trial satisfying \eqref{eq:safe}.
By Lemma~\ref{lem:nonzero}, $r_k>0$.
Taking the inner product of \eqref{eq:kkt-stationarity} with $s_k/r_k$ gives
\begin{equation}\label{eq:linear-gradient}
 \begin{aligned}
 \norm{\nabla f(x_k)+\nabla^2f(x_k)s_k}
 &=
       (2\sqrt{\sigma_kh_k}+\lambda_k)r_k,\\
 (2\sqrt{\sigma_kh_k}+\lambda_k)r_k
 &\le\norm{\nabla f(x_k)}
       +[-\lambda_{\min}(\nabla^2f(x_k))]_+r_k.
 \end{aligned}
\end{equation}
For $0\le u\le1$, \eqref{eq:linear-gradient} implies
\begin{equation}\label{eq:linear-segment}
 \norm{\nabla f(x_k)+u\nabla^2f(x_k)s_k}
 \le\norm{\nabla f(x_k)}
       +[-\lambda_{\min}(\nabla^2f(x_k))]_+r_k.
\end{equation}

By \eqref{eq:model-scales} and \eqref{eq:safe},
\begin{equation}\label{eq:safe-model-scale}
 L_1h_k
 =\max\left\{L_1\norm{\nabla f(x_k)},
       \frac{L_1[-\lambda_{\min}(\nabla^2f(x_k))]_+^2}
            {64\sigma_k}\right\}
 \le\sigma_k.
\end{equation}
Hence,
\begin{equation}\label{eq:safe-radius}
 \sqrt{L_1}r_k
 \le\sqrt{L_1}R_k=
       \frac14\sqrt{\frac{L_1h_k}{\sigma_k}}
 \le\frac14,\qquad L_1r_k^2\le\frac1{16}.
\end{equation}
Let $M_k=\max_{0\le u\le1}\norm{\nabla^3f(x_k+us_k)}$.
Taylor's formula and \eqref{eq:linear-segment} give, for $0\le u\le1$,
\begin{equation}\label{eq:segment-gradient}
 \begin{aligned}
 \norm{\nabla f(x_k+us_k)}
 &\le\norm{\nabla f(x_k)+u\nabla^2f(x_k)s_k}
       +\tfrac12M_kr_k^2\\
 &\le
       \norm{\nabla f(x_k)}
       +[-\lambda_{\min}(\nabla^2f(x_k))]_+r_k
       +\tfrac12M_kr_k^2.
 \end{aligned}
\end{equation}
Taking maxima along the segment gives
\[
 \begin{aligned}
 M_k
 &\underset{\eqref{eq:assumption},\,\eqref{eq:segment-gradient}}{\le}
       L_0+L_1\norm{\nabla f(x_k)}
       +L_1[-\lambda_{\min}(\nabla^2f(x_k))]_+r_k
       +\tfrac12L_1M_kr_k^2\\
 &\underset{\eqref{eq:safe-radius}}{\le}
       L_0+L_1\norm{\nabla f(x_k)}
       +\tfrac14\sqrt{L_1}[-\lambda_{\min}(\nabla^2f(x_k))]_+
       +\tfrac1{32}M_k\\
 &\underset{\eqref{eq:safe}}{\le}\sigma_k+\tfrac1{32}M_k.
 \end{aligned}
\]
Rearranging gives
\begin{equation}\label{eq:third-derivative-bound}
 M_k\le\frac{32}{31}\sigma_k.
\end{equation}
Taylor's formula with \eqref{eq:third-derivative-bound} yields
\begin{align}
 \bigl|f(x_k+s_k)-f(x_k)-\nabla f(x_k)^Ts_k
               -\tfrac12s_k^T\nabla^2f(x_k)s_k\bigr|
 &\le M_kr_k^3/6
       \le\sigma_kr_k^3,
       \label{eq:value-remainder}\\
 \norm{\nabla f(x_k+s_k)-\nabla f(x_k)-\nabla^2f(x_k)s_k}
 &\le M_kr_k^2/2
       \le3\sigma_kr_k^2.
       \label{eq:gradient-remainder}
\end{align}
The optimality conditions \eqref{eq:kkt-stationarity}--\eqref{eq:kkt-psd} give
\begin{equation}\label{eq:model-decrease}
 \begin{aligned}
 -\nabla f(x_k)^Ts_k-\tfrac12s_k^T\nabla^2f(x_k)s_k
 &=
       \tfrac12s_k^T\nabla^2f(x_k)s_k
       +(2\sqrt{\sigma_kh_k}+\lambda_k)r_k^2\\
 &\ge
       \left(\sqrt{\sigma_kh_k}+\tfrac12\lambda_k\right)r_k^2.
 \end{aligned}
\end{equation}
Using $\sigma_kr_k\le\sqrt{\sigma_kh_k}/4$ from
\eqref{eq:model-scales}--\eqref{eq:subproblem},
\begin{align*}
 f(x_k)-f(x_k+s_k)
 &\underset{\eqref{eq:model-decrease},\,\eqref{eq:value-remainder}}{\ge}
       \left(\sqrt{\sigma_kh_k}+\tfrac12\lambda_k\right)r_k^2
          -\sigma_kr_k^3\\
 &\ge\tfrac34\sqrt{\sigma_kh_k}\,r_k^2
          +\tfrac12\lambda_kr_k^2.
\end{align*}
\begin{align*}
 \norm{\nabla f(x_k+s_k)}
 &\underset{\eqref{eq:linear-gradient},\,\eqref{eq:gradient-remainder}}{\le}
       (2\sqrt{\sigma_kh_k}+\lambda_k)r_k+3\sigma_kr_k^2\\
 &\le\tfrac12h_k+\lambda_kr_k+\tfrac3{16}h_k
 =\tfrac{11}{16}h_k+\lambda_kr_k.
\end{align*}
Thus \eqref{eq:accept-value} and \eqref{eq:accept-gradient} hold.
Rejections leave the iterate, and hence the finite threshold in
\eqref{eq:safe}, unchanged while doubling the scale, so
backtracking terminates after finitely many doublings.
\end{proof}

\subsection{Accepted steps: boundary and interior cases}
\label{sec:accepted-steps}

Let $\tau_0<\cdots<\tau_{K-1}$ be the indices of the first
$K\ge1$ accepted trials of Algorithm~\ref{alg:utr}.
Set $\tau_K=\tau_{K-1}+1$, so that $x_{\tau_K}$ is the point
immediately after the $K$-th accepted step.
Rejections leave the point unchanged, so
\begin{equation}\label{eq:accepted-points}
 x_{\tau_0}=x_0,\qquad
 x_{\tau_{j+1}}=x_{\tau_j}+s_{\tau_j}\quad(0\le j<K).
\end{equation}
For these points define the gradient norm, negative curvature,
and function decrease by
\begin{equation}\label{eq:accepted-measures}
 \begin{aligned}
 g_j&=\norm{\nabla f(x_{\tau_j})},\qquad
 \chi_j=[-\lambda_{\min}(\nabla^2f(x_{\tau_j}))]_+
       &&(0\le j\le K),\\
 \Delta_j&=f(x_{\tau_j})-f(x_{\tau_{j+1}})
       &&(0\le j<K).
 \end{aligned}
\end{equation}
Partition the accepted steps into boundary and interior steps:
\begin{equation}\label{eq:step-sets}
 \mathcal B=\{0\le j<K:r_{\tau_j}=R_{\tau_j}\},\qquad
 \mathcal I=\{0,\ldots,K-1\}\setminus\mathcal B.
\end{equation}
With the convention $\sigma_{\tau_{-1}}=2\sigma_0$, the initial
trial for the $j$-th accepted step uses scale $\sigma_{\tau_{j-1}}/2$.
If this trial is accepted, then $\sigma_{\tau_j}=\sigma_{\tau_{j-1}}/2$.
Otherwise, the last rejected trial uses scale $\sigma_{\tau_j}/2$,
which must violate \eqref{eq:safe} by Lemma~\ref{lem:safe}. Hence
\begin{equation}\label{eq:backtrack-upper}
 \sigma_{\tau_j}\le\max\left\{
       \frac{\sigma_{\tau_{j-1}}}{2},\,
       2\bigl(L_0+L_1g_j+\sqrt{L_1}\chi_j\bigr)\right\}
       \quad(0\le j<K).
\end{equation}
By \eqref{eq:accept-value} and $r_{\tau_j}\le R_{\tau_j}$,
\begin{equation}\label{eq:descent-sum}
 \begin{aligned}
 \Delta_j&\ge\tfrac34\sqrt{\sigma_{\tau_j}h_{\tau_j}}\,r_{\tau_j}^2
       \ge3\sigma_{\tau_j}r_{\tau_j}^3>0
                  &&(0\le j<K),\\
 3\sum_{j=0}^{K-1}\sigma_{\tau_j}r_{\tau_j}^3
 &\le\sum_{j=0}^{K-1}\Delta_j
 =f(x_0)-f(x_{\tau_K})\le\Delta.
 \end{aligned}
\end{equation}

\begin{lemma}[Boundary and interior estimates]\label{lem:geometry}
For every $j\in\mathcal B$,
\begin{equation}\label{eq:boundary}
 \begin{aligned}
 g_j&\le16\sigma_{\tau_j}r_{\tau_j}^2,&
 g_{j+1}&\le51\sigma_{\tau_j}r_{\tau_j}^2,\\
 \chi_j&\le32\sigma_{\tau_j}r_{\tau_j},&
 \lambda_{\tau_j}&\le10\sqrt{\sigma_{\tau_j}h_{\tau_j}}.
 \end{aligned}
\end{equation}
For every $j\in\mathcal I$,
\begin{equation}\label{eq:interior}
 h_{\tau_j}=g_j,\qquad \lambda_{\tau_j}=0,\qquad
 g_{j+1}\le\tfrac{11}{16}g_j,\qquad
 \chi_j\le2\sqrt{\sigma_{\tau_j}g_j}.
\end{equation}
\end{lemma}
\begin{proof}
By \eqref{eq:model-scales} and \eqref{eq:accepted-measures},
$g_j\le h_{\tau_j}$ and
$\chi_j\le8\sqrt{\sigma_{\tau_j}h_{\tau_j}}$.
If $j\in\mathcal B$, then $r_{\tau_j}=R_{\tau_j}$, so
\eqref{eq:model-scales} gives
\begin{equation}\label{eq:boundary-identities}
 h_{\tau_j}=16\sigma_{\tau_j}r_{\tau_j}^2,
 \qquad \sqrt{\sigma_{\tau_j}h_{\tau_j}}
       =4\sigma_{\tau_j}r_{\tau_j}.
\end{equation}
Since $r_{\tau_j}=R_{\tau_j}>0$, \eqref{eq:linear-gradient} gives
\begin{equation}\label{eq:boundary-multiplier}
 \begin{aligned}
 \lambda_{\tau_j}
 &\le
       \frac{g_j}{R_{\tau_j}}+\chi_j
       -2\sqrt{\sigma_{\tau_j}h_{\tau_j}}\\
 &\le
       (4+8-2)\sqrt{\sigma_{\tau_j}h_{\tau_j}}
       =10\sqrt{\sigma_{\tau_j}h_{\tau_j}}.
 \end{aligned}
\end{equation}
The gradient test \eqref{eq:accept-gradient} and
\eqref{eq:boundary-identities}--\eqref{eq:boundary-multiplier} give
\[
 \begin{aligned}
 g_{j+1}&\le\tfrac{11}{16}h_{\tau_j}
       +10\sqrt{\sigma_{\tau_j}h_{\tau_j}}\,r_{\tau_j}\\
 &=
       (11+40)\sigma_{\tau_j}r_{\tau_j}^2
       =51\sigma_{\tau_j}r_{\tau_j}^2.
 \end{aligned}
\]

If $j\in\mathcal I$, complementary slackness
\eqref{eq:kkt-complementarity} gives $\lambda_{\tau_j}=0$, and
\eqref{eq:kkt-psd} yields $\chi_j\le2\sqrt{\sigma_{\tau_j}h_{\tau_j}}$.
If $h_{\tau_j}>g_j$, then \eqref{eq:model-scales} implies $\chi_j>0$,
and hence
\[
 \chi_j\le
 2\sqrt{\sigma_{\tau_j}h_{\tau_j}}
 =\frac{\chi_j}{4}<\chi_j,
\]
a contradiction. Thus $h_{\tau_j}=g_j$.
The curvature bound above and \eqref{eq:accept-gradient} give
the remaining estimates.
\end{proof}

\subsection{Scale control and initial gradient reduction}
\label{sec:potential-proof}

\begin{lemma}[One-step scale and gradient estimates]\label{lem:young}
There is a universal constant $C>0$ such that, for every
$0\le j<K$ and $\eta>0$,
\begin{align}
 \sqrt{L_1}\chi_j
 &\le\frac{\sigma_{\tau_j}}{16}+256L_1g_j
                         +CL_1^{3/2}\Delta_j,\label{eq:curvature-young}\\
 L_1g_{j+1}
 &\le\tfrac{11}{16}L_1g_j+\eta\sigma_{\tau_j}
                         +\frac{C}{\sqrt\eta}L_1^{3/2}\Delta_j.
                         \label{eq:gradient-young}
\end{align}
Consequently,
\begin{equation}\label{eq:scale-recursion}
 \sigma_{\tau_j}\le\tfrac47\sigma_{\tau_{j-1}}
                         +3L_0+600L_1g_j+CL_1^{3/2}\Delta_j.
\end{equation}
\end{lemma}
\begin{proof}
If $h_{\tau_j}=g_j$, then
\[
 \sqrt{L_1}\chi_j
       \le
       8\sqrt{\sigma_{\tau_j}L_1g_j}
       \le\frac{\sigma_{\tau_j}}{16}+256L_1g_j.
\]
If $h_{\tau_j}>g_j$, Lemma~\ref{lem:geometry} rules out an interior
step, so $j\in\mathcal B$. By \eqref{eq:model-scales},
$h_{\tau_j}=\chi_j^2/(64\sigma_{\tau_j})$ and
$r_{\tau_j}=R_{\tau_j}=\chi_j/(32\sigma_{\tau_j})$. Hence
\begin{equation}\label{eq:curvature-descent}
 \Delta_j\underset{\eqref{eq:descent-sum}}{\ge}
 3\sigma_{\tau_j}r_{\tau_j}^3
 =\frac{3\chi_j^3}{32^3\sigma_{\tau_j}^2}.
\end{equation}
Young's inequality with exponents $3/2$ and $3$ yields
\[
 \sqrt{L_1}\chi_j
 \underset{\eqref{eq:curvature-descent}}{\le}
 \frac{32}{3^{1/3}}\sigma_{\tau_j}^{2/3}
                      (L_1^{3/2}\Delta_j)^{1/3}
 \le\frac{\sigma_{\tau_j}}{16}+CL_1^{3/2}\Delta_j.
\]
These two cases prove \eqref{eq:curvature-young}.

For $j\in\mathcal I$, \eqref{eq:gradient-young} follows directly
from \eqref{eq:interior}. For $j\in\mathcal B$, Young's inequality gives
\[
 \begin{aligned}
 L_1g_{j+1}
 &\underset{\eqref{eq:boundary}}{\le}51L_1\sigma_{\tau_j}r_{\tau_j}^2\\
 &\le\eta\sigma_{\tau_j}
          +\frac{C}{\sqrt\eta}L_1^{3/2}\sigma_{\tau_j}r_{\tau_j}^3\\
 &\underset{\eqref{eq:descent-sum}}{\le}\eta\sigma_{\tau_j}
          +\frac{C}{\sqrt\eta}L_1^{3/2}\Delta_j.
 \end{aligned}
\]
Together with the interior case, this proves \eqref{eq:gradient-young}.
Moreover,
\[
 \begin{aligned}
 \sigma_{\tau_j}
 &\underset{\eqref{eq:backtrack-upper}}{\le}
 \tfrac12\sigma_{\tau_{j-1}}+2L_0+2L_1g_j+2\sqrt{L_1}\chi_j\\
 &\underset{\eqref{eq:curvature-young}}{\le}
 \tfrac12\sigma_{\tau_{j-1}}+2L_0+514L_1g_j
                  +\tfrac18\sigma_{\tau_j}+CL_1^{3/2}\Delta_j.
 \end{aligned}
\]
Moving $\sigma_{\tau_j}/8$ to the left and rounding up the
coefficients proves \eqref{eq:scale-recursion}.
\end{proof}

\begin{proposition}[Joint potential and scale sums]\label{prop:potential}
Define the potential
\begin{equation}\label{eq:potential-definition}
 P_j=\sigma_{\tau_{j-1}}+10^4L_1g_j\quad(0\le j\le K).
\end{equation}
Then, for a universal constant $C>0$,
\begin{equation}\label{eq:potential}
 P_{j+1}\le\tfrac34P_j+4L_0+CL_1^{3/2}\Delta_j
                  \quad(0\le j<K).
\end{equation}
In particular,
\begin{equation}\label{eq:potential-expansion}
 P_j\le(3/4)^jP_0+16L_0+CL_1^{3/2}\Delta
                  \quad(0\le j\le K),
 \qquad \sigma_{\tau_j}\le CQ_0\quad(0\le j<K).
\end{equation}
For integers $b\ge0$ and $m\ge1$ with $b+m\le K$,
\begin{equation}\label{eq:scale-sum-general}
 \sum_{j=b}^{b+m-1}\sigma_{\tau_j}
 \le3P_b+16mL_0+CL_1^{3/2}\Delta.
\end{equation}
\end{proposition}
\begin{proof}
The weight $10^4$ absorbs the $600L_1g_j$ term in the scale recursion
while retaining contraction. Taking $\eta=10^{-6}$ gives
\[
 \begin{aligned}
 P_{j+1}
 &\underset{\eqref{eq:gradient-young},\,\eqref{eq:potential-definition}}{\le}
       1.01\sigma_{\tau_j}+6875L_1g_j+CL_1^{3/2}\Delta_j\\
 &\underset{\eqref{eq:scale-recursion}}{\le}\tfrac{101}{175}\sigma_{\tau_{j-1}}
          +3.03L_0+7481L_1g_j+CL_1^{3/2}\Delta_j\\
 &\underset{\eqref{eq:potential-definition}}{\le}
       \tfrac34P_j+4L_0+CL_1^{3/2}\Delta_j.
 \end{aligned}
\]
Iterating \eqref{eq:potential} and using \eqref{eq:descent-sum}
proves the bound on $P_j$ in \eqref{eq:potential-expansion}.
Since $P_0=2\sigma_0+10^4L_1g_0\le10^4Q_0$ and
$\sigma_{\tau_j}\le P_{j+1}$, the scale bound follows.

To prove \eqref{eq:scale-sum-general}, sum \eqref{eq:potential}
over $j=b,\ldots,b+m-1$ and rearrange:
\[
 \frac14\sum_{j=b}^{b+m-1}P_{j+1}
 \le\frac34(P_b-P_{b+m})+4mL_0
                         +CL_1^{3/2}\sum_{j=b}^{b+m-1}\Delta_j.
\]
Dropping the nonpositive endpoint term and using
\eqref{eq:potential-definition} and \eqref{eq:descent-sum}
proves \eqref{eq:scale-sum-general}.
\end{proof}

\begin{lemma}[Initial-stage bounds]
\label{lem:initial-stage}
There exists an integer $j_0\ge0$ satisfying
\begin{equation}\label{eq:initial-count}
 j_0=O\left(1+\lp\frac{Q_0}{Q}
              +\lp\frac{g_0}{(\Delta^2Q)^{1/3}}\right)
\end{equation}
such that
\begin{align}
 g_j&\le C(\Delta^2Q)^{1/3}
                         &&(j_0\le j\le K),\label{eq:initial-gradient}\\
 \sum_{j=b}^{b+m-1}\sigma_{\tau_j}
 &\le C\bigl(mL_0+L_1^{3/2}\Delta\bigr)
                         &&(j_0\le b<b+m\le K).
                         \label{eq:scale-sum}
\end{align}
\end{lemma}
\begin{proof}
Set
\[
 j_\sigma=\left\lceil\frac{\lp(P_0/Q)}{\log(4/3)}\right\rceil.
\]
By construction, $(3/4)^{j_\sigma}P_0\le Q$.
Hence,
\begin{equation}\label{eq:scale-stabilized}
 P_j\underset{\eqref{eq:potential-expansion}}{\le}CQ\quad(j_\sigma\le j\le K),\qquad
 \sigma_{\tau_j}\le CQ\quad(j_\sigma\le j<K).
\end{equation}
Also $P_0\le10^4Q_0$ gives
$j_\sigma=O(1+\lp(Q_0/Q))$.
For $j_\sigma\le b<b+m\le K$ with $m\ge1$,
\[
 \begin{aligned}
 \sum_{j=b}^{b+m-1}\sigma_{\tau_j}
 &\underset{\eqref{eq:scale-sum-general},\,\eqref{eq:scale-stabilized}}{\le}
 CQ+16mL_0+CL_1^{3/2}\Delta\\
 &\le C\bigl(mL_0+L_1^{3/2}\Delta\bigr).
 \end{aligned}
\]

For $j\in\mathcal B$,
\begin{equation}\label{eq:boundary-gradient}
 g_{j+1}\underset{\eqref{eq:boundary}}{\le}51\sigma_{\tau_j}r_{\tau_j}^2
       \underset{\eqref{eq:descent-sum}}{\le}
       \frac{51}{3^{2/3}}\sigma_{\tau_j}^{1/3}\Delta_j^{2/3}.
\end{equation}
Induction over the two cases, using $\Delta_j\le\Delta$, gives
\begin{equation}\label{eq:global-gradient-bound}
 g_j\underset{\eqref{eq:interior},\,\eqref{eq:potential-expansion},\,\eqref{eq:boundary-gradient}}{\le}
       \max\{g_0,C(\Delta^2Q_0)^{1/3}\}\qquad(0\le j\le K).
\end{equation}
For $j_\sigma\le j<K$, the sharper scale bound instead gives
\begin{equation}\label{eq:stabilized-gradient-recursion}
 g_{j+1}\underset{\eqref{eq:interior},\,\eqref{eq:scale-stabilized},\,\eqref{eq:boundary-gradient}}{\le}
 \begin{cases}
 C(\Delta^2Q)^{1/3},&j\in\mathcal B,\\
 (11/16)g_j,&j\in\mathcal I.
 \end{cases}
\end{equation}
Thus, boundary steps bound the gradient at the target scale, while
interior steps reduce it geometrically.
With a common universal $C\ge1$, iterating
\eqref{eq:stabilized-gradient-recursion} from \eqref{eq:global-gradient-bound}
gives, for $j_\sigma\le j\le K$,
\[
 g_j\le\max\left\{
 (11/16)^{j-j_\sigma}
       \max\{g_0,C(\Delta^2Q_0)^{1/3}\},\,
 C(\Delta^2Q)^{1/3}\right\}.
\]
Since $C\ge1$ and $Q_0\ge Q$,
\[
 \frac{\max\{g_0,C(\Delta^2Q_0)^{1/3}\}}{C(\Delta^2Q)^{1/3}}
 \le\max\left\{1,\frac{g_0}{(\Delta^2Q)^{1/3}}\right\}
       \left(\frac{Q_0}{Q}\right)^{1/3}.
\]
It therefore suffices to take
\begin{equation}\label{eq:initial-index}
 j_0=j_\sigma+
 \left\lceil
 \frac{\displaystyle\lp\frac{g_0}{(\Delta^2Q)^{1/3}}
             +\frac13\log\frac{Q_0}{Q}}
      {\log(16/11)}
 \right\rceil.
\end{equation}
The choice \eqref{eq:initial-index} bounds the first term in the maximum
by $C(\Delta^2Q)^{1/3}$, proving
\eqref{eq:initial-count}--\eqref{eq:initial-gradient}.
Since $j_0\ge j_\sigma$, \eqref{eq:scale-sum} also follows.
\end{proof}

